\newpage

\insimg{forprint/134_folio_39_recto_B_OF_KP_II.jpg}
\vskip -10pt

\Rline{\krnt, folio 39\,\RE.}
\bigskip

\begin{sans}
koraṇṭakapaddhatau vajrolīmudrā\\ % Main Heading centered, bold
\smallskip

\centerline{śrīgaṇeśāya namaḥ}% salutation centered
\smallskip

yasya syurvītasaṃkalpāḥ prāṇendriyamanodhiyām\!।\\ %linebreak
vṛttayastasya\endnote{{vṛttayas tasya~] \textit{diagnostic conj.}, tasya tasyaiva M. Cf. \textit{Bhāgavatapurāṇa} 11.11.14 (yasya syur vītasaṅkalpāḥ prāṇendriyamanodhiyām~| vṛttayaḥ sa vinirmukto dehastho 'pi hi tadguṇaiḥ). The verse in the \textit{koraṇṭakapaddhati} has been adapted from \textit{Bhāgavatapurāṇa} 11.11.14. It is likely that the word \textit{vṛttayaḥ} has dropped out of the \textit{Koraṇṭakapaddhati}’s opening verse because without it the syntax does not make sense.} mudrā syānnānyo\endnote{nānyo~] \textit{conj.}, nonyo M.} bhavitumarhati\!॥1\!॥\\ %linebreak
sahāyasya tu puṇyaṃ syāddhanavṛddhirbhaviṣyati\!।\\ %linebreak
kṛcchracāndrāyaṇaiḥ\endnote{kṛcchra°~] \textit{corr.}, kṛchra° M.} puṇyaṃ\endnote{puṇyaṃ~] \textit{conj.} \textsc{Goodall}, puṇyair M.} mudrayā bhavati dhruvam\!॥2\!॥\\ %linebreak
puṇyenānena svargaṃ syādātmalabdhyā\endnote{°labdhyā~] \textit{conj.} \textsc{Goodall}, °labdhvā M.} sukhī bhavet\!।\\%linebreak
viṣṇornirmālya vajrolī śivanirmālya maithunam\!॥3\!॥\\ %linebreak
tasmātsarvaprayatnena\endnote{°yatnena~] \textit{emend.}, °yatne M.} vajrolīṃ\endnote{vajrolīṃ~] \textit{corr.}, vajrolī M.} sādhayetsudhīḥ\!।\\%linebreak
vajrolīṃ śaṃsayiṣyāmi hyūrdhvabindorvidhāraṇe\!॥4\!॥\\ %linebreak
bindunā bhidyate nāḍī vajrakandarasaṃjñakā\endnote{°kandara°~] \textit{corr.}, °kaṃdhara° M. °saṃjñakā~] \textit{emend}., °saṃjñake M.}\!।\\%linebreak
tasmādvajrolī vikhyātā matsyendreṇa prakāśitā\!॥5\!॥\\ %linebreak
prathamaṃ vṛṣapādādisaṃsthā āsanamucyate\!।\\%linebreak
abhyāsenaiva tatrādau\endnote{tatrādau~] \textit{conj.}, tatrāddhā M.} āsaneṣu vinigrahaḥ\!।\\%linebreak
abhyāsenāsane siddhe gudakambukriyā smṛtā\!॥6\!॥\\ %linebreak
yathā yathā lapsyati\endnote{lapsyati~] \textit{emend.}, lapsyasi M.} kaṣṭameṣaḥ\\ %linebreak
\hspace*{1cm}tathā tathā śudhyati nāḍicakram\!।\\%linebreak
ātmopalabdhyātiviśuddhatattvaḥ\endnote{°labdhyā°~] \textit{emend}. \textsc{Goodall}, °labdhvā° M.}\\ %linebreak
\hspace*{1cm}saṃsāracakraiḥ parimucyate\endnote{parimucyate~] \textit{conj.}, parimadhyate M.} 'sau\endnote{'sau~] \textit{avagraha added}, sau M.}\!॥7\!॥\endnote{The metre of verse 7 is \textit{upajāti}.}\\ %linebreak
baḍavākambuvatkāryamākuñcanaprasāraṇam\endnote{°kambu°~] \textit{conj.}, °kaṃcu° M. Cf. \textit{Kuraṇṭakapaddhati}: naulisiddhyartham aśvādhārakambuvad vāraṃ vāraṃ kambum ākuñcayet|}\!।\\%linebreak
dvihastayostarjanībhyāṃ kāryaṃ\endnote{kāryaṃ~] \textit{diagnostic conj.}, savya° M.} savyāpasavyataḥ\!।\\%linebreak
pradakṣiṇavadgharṣaṇaṃ syāditi kambukriyā matā\!॥8\!॥\\%linebreak
dvisahasratrisahasraṃ sahasraṃ ca caturthakam\!।\\%linebreak %\endnote{The first verse quarter is unmetrical. Strictly speaking, the sixth syllable should be long.}
evaṃ ṣoḍaśaparyantaṃ kartavyaṃ ca dine dine\!॥9\!॥\\%linebreak
prathamābhyāsaprārambhe lagati praharadvayam\!।\\%linebreak
evaṃ ṣoḍaśasāhasramabhyāsastu muhūrtataḥ\!॥10\!॥\\%linebreak
\end{sans}
\medskip

\begin{eng}
\titleinline{\color{deepmaroon}{\bfseries [Koraṇṭaka’s Manual on Vajrolīmudrā]}}

\centerline{{\large Homage to Gaṇeśa.}}
\end{eng}

\begin{multicols}{2}
\begin{eng}
\heading{[Introduction]}
\begin{enumerate}
\item[1.] One whose actions of breath, senses, mind, and intellect are free from desires may attain [\tit{vajrolī}]\tit{mudrā}. No other person is worthy.
\item[2.] By means of this \tit{mudrā}, [even] one’s friend attains merit, and their wealth will increase. The merit acquired by the Arduous and Lunar fasts (\tit{kṛcchracāndrāyaṇa}) is sure to arise.
\item[3.] Through the merit acquired, heaven is gained, and by attaining the Self, one becomes blissful. That which belongs to Viṣṇu is \tit{vajrolī}. That which belongs to Śiva is [mundane] sex.
\item[4.] Therefore, the wise person should practise \tit{vajrolī} with all effort. I praise \tit{vajrolī} for retaining upward-flowing semen (\tit{ūrdhvabindu}).
\item[5.] The channel called the “hollow of the penis” (\tit{vajrakandara}) is pierced by semen. Therefore, Matsyendra named \tit{vajrolī} [thus] and brought it to light.
\end{enumerate}

\heading{[Preliminary Practices for \tit{Vajrolī}]}

\heading{[Posture (\textit{āsana})]}

\begin{enumerate}
\item[6.] First of all, \tit{āsana} is taught as the arrangement of [postures that] begin with [Pawing] the Leg like
a Bull (\tit{vṛṣapāda}). In this regard, holding the postures [steadily] first arises through practice. When
\tit{āsana} has been perfected through practice, the technique of \tbf{[Cleansing] the Anus} is
recommended.
\item[7.] The more [the yogi] experiences pain, the more his network of channels is purified. By perceiving
the Self, he becomes one for whom all the ontic principles (\tit{tattva}) have been completely purified
and is freed from the cycles of transmigration.
\end{enumerate}

\heading{[The Technique of Cleansing the Anus (\tit{gudakambukriyā})]}

\begin{enumerate}
\item[8.] One should contract and expand [the anus] like the anus of a mare. One should rub [the anus] with
the index fingers of both hands, [circling] to the left and right like a circumambulation. This is
regarded as the technique of \tbf{[Cleansing] the Anus} (\tit{kambukriyā}).

\item[9.] One should do it two, three, and four thousand times, and thus up to sixteen 
[thousand] times per day.

\item[10.] Initially, upon undertaking the practice, it takes six hours. In this way, [after some time,] sixteen
thousand [repetitions are done] in forty-eight minutes (\tit{muhūrta}).
\end{enumerate}

\end{eng}
\end{multicols}
\bigskip


%~ \begin{multicols}{2}
%~ \begin{kan}
%~ \end{kan}
%~ \end{multicols}

\newpage

\insimg{forprint/135_folio_39_verso.jpg}
\vskip -10pt

\Rline{\krnt, folio 39\,\VE.}
\bigskip

\begin{sans}
yadā\endnote{yadā~] M(pc), yathā M(ac).} bhavati yogeśastadā vai bhavati dhruvam\!।\\%linebreak
kambukriyānantaraṃ vai bastiḥ kāryā vicakṣaṇaiḥ\!॥11\!॥\\%linebreak
prādeśamātrā dīrghe sā sthūlāṅguṣṭhapramāṇataḥ\!।\\%linebreak
vijñeyā nalikā bastau kriyā kāryāvicārataḥ\!॥12\!॥\\%linebreak
nābhidaghne jale nyasya pāyau nālīṃ samarpayet\!।\\%linebreak
sahasrakṛtva\endnote{°kṛtva~] \textit{corr.} \textsc{Goodall}, °kṛtvam M.} udakamudaraṃ pratipūritam\!।\\%linebreak
tyaktvā tyaktvā tu gṛhṇīyādevaṃ bastikriyā bhavet\!॥13\!॥\\%linebreak
nalikāyā ṛte vāyurjalaṃ pāyau praveśya vai\!।\\%linebreak
jīryate hi tadā bastiḥ siddhā\endnote{bastiḥ siddhā~] \textit{emend.} \textsc{Mallinson}., bastisiddha M.} vakti sma kanthaḍī\!॥14\!॥\\%linebreak
varāhamāsanaṃ kṛtvā vāyumākṛṣya pāyunā\!।\\%linebreak
bhūya āmucya vāyuṃ tu śataśo 'tha\endnote{'tha~] \textit{avagraha added}, tha M.} sahasraśaḥ\!॥15\!॥\\%linebreak
bhastrākhyo\endnote{bhastrā°~] \textit{emend.}, bhasrā° M.} vāyuprayogaḥ\!। kuṇḍalīdvayaṃ prabodhayati\!। $\dagger$janayatyānāṃ$\dagger$ viśeṣaṃ nābhiṃ bhittvā\endnote{bhittvā~] \textit{corr.}, bhitvā M.} bhinatti hṛdayam\!।\\%linebreak prose passage (so no verse number)
liṅgaśuddhiṃ tataḥ kuryātsadgurūttamaśikṣayā\!।\\%linebreak
mūtraṃ nirodhayelliṅge yāvadvedanayā sahet\!॥16\!॥\\%linebreak
tato mucyedbindumātraṃ punaḥ saṃrodhayecchanaiḥ\endnote{°rodhayec~] \textit{corr.}, °rodhaye M.}\!।\\%linebreak
saṃniruddhaṃ punarmucyetpunarmucyetpunaḥ punaḥ\!॥17\!॥\\%linebreak
\end{sans}
\medskip

\begin{multicols}{2}
\begin{eng}
\begin{enumerate}
\item[11.] When one becomes a master of yoga (\tit{yogeśa}), this will definitely be the case. After [practising] the
technique of \tbf{[Cleansing] the Anus}, adepts should practise the \tbf{Enema}.
\end{enumerate}

\heading{[The Enema (\tit{basti})]}

\begin{enumerate}
\item[12.] The tube (nālikā) for the enema is known to be twelve finger-breaths in length and its thickness is
the measure of a thumb’s width. The technique should be done without hesitation.
\item[13.] After placing [oneself ] in water up the navel, one should insert the tube into the anus. A thousand
times, the stomach is filled with water. [The water] is released repeatedly and drawn in. Thus the
Enema Technique is [done].
\item[14.] When one makes water enter the anus without the tube, wind [in the abdomen] is digested. Then,
the Enema [Technique] is perfected. Kanthaḍī taught this.
\item[15.] One should do the Boar Pose, draw in air through the anus, and release it again, [doing thus]
hundreds or thousands of times.
\end{enumerate}

This breath practice is called the Bellows (\tit{bhastrā}). It awakens the two \tit{kuṇḍalinīs}. It generates the
distinct †unstruck sound†. Having pierced the navel, \tit{kuṇḍalinī} pierces the heart.

\heading{[Purifying the Penis (\tit{liṅgaśuddhi})]}

\begin{enumerate}
\item[16.] Now, one should purify the penis according to the teaching of an excellent, true guru. One should
stop the [flow of ] urine in the penis for as long as one can bear with the pain.
\item[17.] Then one should release only a drop and gradually stop [the flow] again. Then the held [urine]
should be released [a little] and released again and again.
\end{enumerate}
\end{eng}
\end{multicols}
\bigskip


%~ \begin{multicols}{2}
%~ \begin{kan}
%~ \end{kan}
%~ \end{multicols}

\newpage

\insimg{forprint/136_folio_40_recto.jpg}
\vskip -10pt

\Rline{\krnt, folio 40\,\RE.}
\bigskip

\begin{sans}
jātyaṅkurasamā sthūlā mūle mṛdutamā matā\!।\\%linebreak
hastapramāṇā dīrghe tu hariccharaśalākikā\endnote{haric°~] \textit{corr.}, hari° M.}\!।\\%linebreak
liṅgachidre praveṣṭavyā śanaiḥ śīghraṃ śanaiḥ śanaiḥ\!॥18\!॥\\%linebreak
prathamaṃ cāmṛtalatā praveṣṭavyā prayatnataḥ\!।\\%linebreak
tatpraveśe vedanayā jvaro bhavati vai mahān\!॥19\!॥\\%linebreak
tannivṛttiśca\endnote{°vṛttiś~] \textit{emend.}, °vṛttiṃ M.} tatroktā strīguṇānāṃ ca varṇanam\!।\\%linebreak
nistejaske varṇane tu prekṣaṇaṃ bahusaṃmatam\endnote{°mataṃ~] \textit{emend.}, °mitaṃ M.}\!॥20\!॥\\%linebreak
ālokane nistejaske\endnote{nistejaske~] \textit{corr.}, nitejaske M.} sāmīpye keli saṃmatā\endnote{Correcting \textit{keli} to \textit{keliḥ} would be unmetrical.}\!।\\%linebreak
sāmīpyakelirvyarthā ca yadā bhavati yoginaḥ\endnote{yoginaḥ~] \textit{emend.}, yoginā M.}\!।\\%linebreak
tadā manyeta vai yogī yuvatīpāṇisparśanam\!॥21\!॥\\%linebreak
śithile pāṇisparśe tu stanasparśo yathāyatham\!।\\%linebreak
vyarthībhūte tu ca sparśe mukhacumbanamiṣyate\!॥22\!॥\\%linebreak
gaṇḍaṃ gaṇḍena saṃghṛṣya nitambaṃ gharṣayennakhaiḥ\!।\\%linebreak
sakāmāṃ yuvatīṃ kṛtvā manāṅmaithunakelivat\endnote{°keli°~] \textit{conj.}, °keki° M.}\!॥23\!॥\\%linebreak
bindau kṣubdhe nalikayā liṅge vāyuṃ praveśayet\!।\\%linebreak
nalyāṃ vāyuṃ pragṛhṇīta\endnote{°gṛhṇīta~] \textit{corr.}, °gṛṇhīta M.} udakaṃ\endnote{udakaṃ~] \textit{conj.} \textsc{Mallinson}, codakaṃ M. It is likely that \textit{ca} was inserted later to separate the vowels as it is a major defect to begin a verse quarter with \textit{ca}.} pāradaṃ\endnote{pāradaṃ~] \textit{corr}., pāradhaṃ M.} tathā\!॥24\!॥\\%linebreak
dugdhaṃ tailaṃ tathā snigdhaṃ picumandaṃ ca rañjanī\endnote{rañjanī~] \textit{corr.}, rājanī M.}\!।\\%linebreak
lavaṅgaṃ jātipattraṃ\endnote{°pattraṃ~] \textit{corr.}, patraṃ M.} ca karpūraṃ kuṅkumaṃ tathā\!॥25\!॥\\%linebreak
kastūrīṃ candanaṃ bhadrāṃ brāhmīṃ\endnote{ brāhmīṃ~] \textit{corr.},  brāhmī°} mocarasaṃ tathā\!।\\%linebreak
etatsuvāsitaṃ\endnote{°vāsitaṃ~] \textit{emend.},  °vāsitu M.} grāhyaṃ jalaṃ\endnote{jalaṃ~] \textit{corr.}, jala° M.} nālena yoginā\endnote{yoginā~] \textit{conj.}, yogavit M.}\!॥26\!॥\\%linebreak
nalikā vakrā kuṭilā saralā viralā $\dagger$tanotarā$\dagger$ mṛdulā\!।\\%linebreak
gurudīkṣitena kāryā\endnote{kāryā~] \textit{emend}. \textsc{Goodall}, kuryād M.} dhāryā sarvaiva\endnote{sarvaiva~] \textit{conj.} \textsc{Goodall}, sarvadaiva M. This conjecture restores the metre, which is \textit{gīti}.} liṅganāle na\!॥27\!॥\\%linebreak
caturaṅguladīrghā ṣaḍaṅgulāṣṭādaśāṅgulā\!।\\%linebreak
tato 'dhikāṅgulā\endnote{'dhikā°~] \textit{avagraha added}, dhikā° M. °ṅgulā~] \textit{emend}., °ṅgulāś M.} ca caturviṃśatiparyantā $\dagger$ [...]$\dagger$ \!॥28\!॥\\%linebreak
\end{sans}

\medskip

\begin{multicols}{2}
\begin{eng}
\begin{enumerate}
\item[18.] A very soft, small probe of \tit{haritaśara} grass, a cubit in length and at its base as thick as a jasmine
sprout, should be inserted into the opening of the penis, gradually, [then] quickly, [then] very
gradually.
\item[19.] But first, one should carefully insert a stalk of the \tit{amṛtā} vine. When it has been inserted, an intense
fever arises along with pain.
\end{enumerate}

\heading{[The Practice of Vajrolī]}

\heading{[Amorous Play with a Woman]}

\begin{enumerate}
\item[20.] [In Kuraṇṭaka’s tradition,] suppressing the fever is prescribed. And [then the yogi recites] a
description of the beautiful qualities of a woman. When the description does not stimulate [him],
looking [at a woman] is highly recommended.

\item[21.] When looking at her does not stimulate him, amorous play near her is allowed. When amorous play
in proximity [to her] has no effect on the yogi, the yogi should think of touching the hand of a
young woman.

\item[22.] When touching her hand is dull, he should touch her breast in due sequence. When touching her
becomes pointless, kissing her mouth is prescribed.

\item[23.] Having rubbed her cheek with his cheek, he should rub her buttocks with his fingernails. [Then,]
he should make her lustful by doing a little love-play.

\item[24.] If his semen stirs, he should make air enter his penis through the tube. He should draw in air
through the tube [and], in the same way, water and mercury.
\end{enumerate}

\heading{[Herbal Urethral Enemas]}

\begin{enumerate}
\item[25--26.] Through the tube, the yogi should take in water that has been deeply infused with the
following: milk, sesame oil, the red castor-oil plant, neem, turmeric, cloves, jasmine leaves,
camphor, saffron, musk, sandal, sweet flag, brāhmī, and the resin of the silk cotton tree.

\item[27.] One initiated by a guru should make a small tube that is crooked, curved, straight, fine, †very
slender† or soft. The whole [tube] should not be put in the urethra.

\item[28.] [The tube’s] length is four, six, eighteen or more finger-breadths, up to twenty-four.
\end{enumerate}

\end{eng}
\end{multicols}
\bigskip


%~ \begin{multicols}{2}
%~ \begin{kan}
%~ \end{kan}
%~ \end{multicols}


\newpage

\insimg{forprint/137_folio_40_verso.jpg}
\vskip -10pt

\Rline{\krnt, folio 40\,\VE.}
\bigskip

\begin{sans}
$\dagger$aṅgulībhiḥ guruśikṣayābhyasitā$\dagger$\endnote{It appears that some syllables and words have dropped out of 28cd and 29ab as 28ab and 29cd are metrical (\textit{anuṣṭubh}). Also, 29ab does not make sense.}\!।\\%linebreak
śatādhikā nirmitavyā nalikā guruśikṣayā\!॥29\!॥\\%linebreak 
vāsudevasya pūjā hi kartavyā vighnaśāntaye\!।\\%linebreak
japedupāṃśunā yogī sarvāvasthāsu buddhimān\!॥30\!॥\\%linebreak 
abhyāse madhyame 'siddhe\endnote{'siddhe~] \textit{avagraha added}, siddhe M.} mantro dvādaśavarṇavān\!।\\%linebreak
oṃkārarahito jāpyo vidhīyeta prayatnataḥ\!॥31\!॥\\%linebreak 
samagrābhyāsasiddhau ca manumoṃkāramāditaḥ\!।\\%linebreak
sadā sarvatra hi japeddeśakāle ca yogavān\!॥32\!॥\\%linebreak 
yamāḥ\endnote{yamāḥ~] \textit{corr.}, yamās M.} sarve yāmitavyāḥ\endnote{The verb form \textit{yāmitavyāḥ} is unattested. The correct form is \textit{yamitavyāḥ} but this would be unmetrical.} niyamāśca kvacitkvacit\!।\\%linebreak
vāsudevasya pūjā tu kartavyā bahubhaktitaḥ\!।\\%linebreak
vāsudevasya mantreṇa\endnote{mantreṇa~] M(pc), ++yana M(ac).} dvādaśākṣaravidyayā\!॥33\!॥\\%linebreak 
aṣṭāṅgaṃ maithunaṃ kāryaṃ pradhānāṅgena varjitam\!।\\%linebreak
pradhānāṅgaṃ bindupātaḥ śeṣamaṣṭāṅgamaithunam\!॥34\!॥\\%linebreak
śiśnena rasamaśnīyāj jihvāmaithunamācaret\endnote{jihvā~] \textit{emend}., jihvāṃ M.}\!।\\%linebreak
maithunaṃ kuhare jihvāpraveśastu prayatnataḥ\!॥35\!॥\\%linebreak
prāśanaṃ raja{}ādānaṃ sahajolīti saṃjñitam\!।\\%linebreak
amarolī prayatnena kartavyā yogamicchatā\endnote{icchatā~] \textit{corr.}, ichatā M.}\!।\\%linebreak
amarīṃ\endnote{amarīṃ~] \textit{emend.}, amarī M.} ca pibedvidvāntsarvarogasya\endnote{On the unconventional \textit{sandhi} of °\textit{nt s}°, see Oberlies 2003: 50. We wish to thank Dominic Goodall for the reference.} śāntaye\!॥36\!॥\\%linebreak
iti nalikākarṣaṇakriyābhirāśleṣopabhoganarmabhiśca\!।\\%linebreak
puṃsāṃ sidhyati $\dagger$sidhyati sidhyati dūrā triṃśati viṃśati vartya\endnote{vartya~] M(pc), varṣa M(ac).} māsaiḥ$\dagger$\!॥37\!॥\\%linebreak
\end{sans}
\medskip

\begin{multicols}{2}
\begin{eng}
\begin{enumerate}
\item[29.] †The practice with a tube of these lengths should be done as taught by the guru.† More than one
hundred [different] tubes should be made according the guru’s teaching.
\end{enumerate}

\heading{[Mantra Recitation]}

\begin{enumerate}
\item[30.] The worship of Vāsudeva should be done to eliminate obstacles. In a low voice, the wise yogi should
repeat [Vāsudeva’s mantra] in all conditions.
\item[31.] When an intermediate level of practice is not [yet] accomplished, the twelve-syllable [Vāsudeva]
mantra should be recited without the [initial] OṂ. One should observe this practice carefully.
\item[32.] Once the entire practice has been accomplished, one should recite the mantra with OṂ at the
beginning. Always and everywhere, indeed, the one who [truly] possesses yoga should recite it in
every place and time.
\item[33.] All the behavioural codes (\tit{yama}) should be upheld, and the observances (\tit{niyama}) wherever
appropriate. As for the worship of Vāsudeva, it should be done with great devotion using the
mantra of Vāsudeva, [that is,] the mantra with twelve syllables.
\end{enumerate}

\heading{[\tit{Sahajolī} and \tit{Amarolī}]}

\begin{enumerate}
\item[34.] Sex with eight components should be performed without the primary component. The primary
component is the emission of semen. The other components [should still be practised].
\item[35.] [The yogi] should consume sexual fluids (\tit{rasa}) through his penis. He should [also] practice tongue
sex. [Tongue] sex is carefully entering the tongue into the vagina.
\item[36.] Eating [in the sense of ] taking in female sexual fluid (\tit{rajas}) is known as sahajolī. Amarolī should be
done carefully by one desirous of yoga. And the wise person should drink their own urine (\tit{amarī})
to cure all diseases.
\item[37.] Thus, men achieve [\tit{sahajolī} and \tit{amarolī}] through the techniques of drawing in [fluids] via the tube
and through the amorous play of enjoying intimate contact [with women]. †The practice is
definitely accomplished in twenty or thirty months.†
\end{enumerate}
\end{eng}
\end{multicols}
\bigskip


%~ \begin{multicols}{2}
%~ \begin{kan}
%~ \end{kan}
%~ \end{multicols}

\newpage

\insimg{forprint/138_folio_41_recto.jpg}
\vskip -10pt

\Rline{\krnt, folio 41\,\RE.}
\bigskip

\begin{sans}
ūrdhvabindostu mahimā na bhūto na bhaviṣyati\!।\\%linebreak
urdhvabindorṛte\endnote{urdhva~] \textit{emend}., urdhvaṃ M.} jñānaṃ hṛdaye na ca tiṣṭhati\!॥38\!॥\\%linebreak
jñānaṃ vinā hyātmalābhaṃ na ca vindati yogavit\!।\\%linebreak
ātmalābhaṃ vinā mṛtyuṃ nāyaṃ tarati yogavān\!॥39\!॥\\%linebreak
tasmātsarvaprayatnena yatnaḥ karyaḥ prayatninā\!।\\%linebreak
katibhirjanmabhiḥ\endnote{janmabhiḥ~] \textit{corr.}, janmabhis M.} siddhirbhaviṣyati na saṃśayaḥ\!॥40\!॥\\%linebreak
tatra ṣoḍaśabhiḥ\endnote{ṣoḍaśabhiḥ~] \textit{corr.}, ṣoḍaśabhi M.} strībhiḥ ratiḥ syānnirbharā yadā\!।\\%linebreak
bindudhāraṇasāmarthyaṃ śoṣaṇaṃ gilanaṃ bilāt\!।\\%linebreak
tadā bindujayo jñeyo yoginā sukhamicchatā\!॥41\!॥\\%linebreak
gauraṃ ślakṣṇaṃ suvinītaṃ duṣṭarogādivarjitam\!।\\%linebreak
tāruṇyāvayavaiḥ śuddhaṃ nyūnādhikyavivarjitam\!॥42\!॥\\%linebreak
gāyakaṃ kaṇṭhamadhuraṃ tantrīvādanadakṣiṇam\!।\\%linebreak
rāmāṣoḍaśakaṃ jñeyaṃ vajrolyāṃ bahusaṃmatam\!॥43\!॥\\%linebreak
ekaikāmāśliṣanbhuktvā\endnote{{°śliṣan~] \textit{conj.} \textsc{Goodall}, °śliṣāṃ M.} nalikāṃ ca punaḥ punaḥ\!।\\%linebreak
prārambhe rogajālāni vighnajālaṃ tathaiva ca\!।\\%linebreak
dambhāya\endnote{dambhāya~] \textit{conj.} \textsc{Goodall}, dambhāyoś M.} cātha\endnote{cātha~] \textit{conj.}, catha M.} vighnā ye prabhavantyeva yoginaḥ\!॥44\!॥\\%linebreak
ārabdhāyāṃ tu mudrāyāṃ vedanā vāyukarṣaṇāt\!।\\%linebreak
gharṣaṇātsarvanāḍīnāṃ cakrādīnāṃ vyatikramāt\!॥45\!॥\\%linebreak
prabhavatyavisahyā sā sahyā tairjapanāmabhiḥ\endnote{tair~] \textit{emend.},  °tsā M.}\!।\\%linebreak
keśavādīni nāmāni kṛṣṇāntāni kṛtāni yat\!।\\%linebreak
japastu vāsudeveti kavibhiḥ parikīrtitaḥ\!॥46\!॥\\%linebreak
abhyāse madhyame jāte kaṇḍūtirvedanā bhavet\!।\\%linebreak
abhyāse hi dṛḍhībhūte kaṇḍūtirharṣatāṃ vrajet\!॥47\!॥\\%linebreak
\end{sans}
\medskip

\begin{multicols}{2}
\begin{eng}
\heading{[Upward-flowing Semen]}
\begin{enumerate}
\item[38.] There has never been, nor will there ever be, a greatness [surpassing] upward-flowing semen. And
without upward-flowing semen, gnosis does not endure in the heart.
\item[39.] Without gnosis, the adept of yoga does not attain the Self. Without the attainment of the Self, a
yogi does not overcome death.
\item[40.] Therefore, with utmost effort, the diligent person should strive. Success will arise within a few
lifetimes. There is no doubt.
\item[41.] In this [yoga], when a yogi who desires bliss has passionate sex with sixteen women and can retain
semen, drying up and swallowing [fluids] from the vagina, he then knows mastery of semen.
\item[42--43.] One should understand that the sixteen beautiful women highly recommended for \tit{vajrolī} should
be fair-skinned, soft, well-behaved, free from serious diseases, pure, with youthful limbs, and
without missing or extra body-parts. They should be singers with sweet voices and skilled in
playing stringed instruments.
\item[44.] Embracing each one, he should enjoy her, repeatedly [inserting] the tube [into his urethra]. In the
beginning, multitudes of diseases and a multitude of obstacles arise. Now, those obstacles that arise
in order to delude him overpower the yogi.
\item[45.] When [\tit{vajrolī}]\tit{mudrā} is [first] undertaken, pain arises due to drawing in the air, its friction on all of
the channels, and the disturbance of \tit{cakras} and the like.
\item[46.] [If ] unbearable pain overwhelms him, it becomes bearable by reciting the [twenty-four] names [of
Vāsudeva], which begin with Keśava and end with Kṛṣṇa. Sages praise the recitation of VĀSUDEVA.
\item[47.] When an intermediate [stage of ] practice arises, the pain becomes an itching [sensation]. When the
practice becomes well-established, the itching turns into joy.
\end{eng}
\end{multicols}
\bigskip


%~ \begin{multicols}{2}
%~ \begin{kan}
%~ \end{kan}
%~ \end{multicols}

\newpage

\insimg{forprint/139_folio_41_verso_E_OF_KP_II.jpg}
\vskip -10pt

\Rline{\krnt, folio 41\,\VE.}
\bigskip

\begin{sans}
abhyāsasaṅkhyā kartavyā keśavādyabhidhānataḥ\!।\\%linebreak
ekādibhirna kartavyā hyaṅkaiḥ\endnote{hyaṅkaiḥ~] \textit{conj}, vyaṅkaiḥ M.} prākṛtasaṃskṛtaiḥ\!।\\%linebreak
nāmabhirjāyate puṇyam aṅkaiśca vitathībhavet\!॥48\!॥\endnote{Cf. \textit{Kuraṇṭakapaddhati}: vāsudeveti manasā japet| abhyāsasaṅkhyāṃ keśavādicaturviṃśatināmabhiḥ kuryāt| puṇyasañcayo bhavati|}\\%linebreak
liṅgachidre nalīṃ dadyād āsyena dhamanaṃ bahu\!।\\%linebreak
kuryādgṛhīta vāyuṃ vai vāyorgrahaṇamucyate\!॥49\!॥\\%linebreak
muṣkasyādhaḥ sīvanīṃ kuñcayitvā muktvā kuryāccālanaṃ\endnote{kuryāc~\textit{conj.}, kuñcec M.} nirbharaṃ ca\!।\\%linebreak
vāraṃ vāraṃ\endnote{vāraṃ vāraṃ~] \textit{conj.}, cālaṃ cālaṃ M.} cālayetkārayedvā $\dagger$kāraṃ kāramutkālabalaṃ ca$\dagger$\!॥50\!॥\endnote{Cf. \textit{Kuraṇṭakapaddhati}: bāhyavāyuṃ gṛhṇiyāt tyajet| evaṃ vāraṃ vāraṃ sahasradvisahasratrisahasram abhyāsaṃ kuryāt| vṛṣaṇasyādhaḥ vāraṃ vāram ākuñcayet| saraṭakabalaṃ bhavati|}\\%linebreak
nalinyā vāyu\endnote{nalinyā~] \textit{emend.}, nalikā° M.} saṃgṛhya tarjanyaṅguṣṭhakartarīm\endnote{°aṅguṣṭha~] \textit{corr.}, °āṅguṣṭha M. kartarīm~] \textit{corr.}, kartarī M.}\!।\\%linebreak
kṛtvā dhṛtvā maṇiṃ liṅge huṃhuṃkāreṇa karṣayet\!॥51\!॥\\%linebreak
visṛjec cāryakāreṇa phūtkāreṇa\endnote{phūt°~] \textit{conj.}, ghāt° M} ca vai punaḥ\!।\\%linebreak
huṃhuṃ kṛtvā tu phūtkāraṃ cāryakāreṇa vai punaḥ\!॥52\!॥\endnote{52cd is an insertion from the lower margin.}\\%linebreak
kapotakaṇṭhamabhyāsaṃ gahano vakti vai khalu\!।\endnote{Cf.\textit{Kuraṇṭakapaddhati}: saraṭakabalaṃ kṛtvā bastipradeśam ūrdhvam ānayet| pārāvatakaṇṭhaṃ bhavati|}\\%linebreak
maṇḍūkodaramabhyāsaḥ kartavyo yoginā śanaiḥ\!॥53\!॥\\%linebreak
nalinyā vāyuṃ\endnote{vāyuṃ~] \textit{corr.}, vāyu M. This restores the metre to Upagīti.} saṃgṛhya bastiṃ\endnote{bastiṃ~] \textit{corr.}, basti M.} saṃpūrya yuktitaḥ\!।\\%linebreak
maṇḍūkodaravatkāryaṃ bhekatundo 'yamucyate\!॥54\!॥\\%linebreak
bhekodarādantaraṃ\endnote{bhekodarād~] \textit{conj.}, bhekaudarān M.} vai caṣakābhyāsa ucyate\!।\\%linebreak
liṅgadvāragṛhīto yo nalikā vāyurantike\!।\\%linebreak
tasyākarṣaṃ\endnote{°karṣaṃ~] \textit{emend.}, °karṣo M.} nābhidvāre kuryādvai haṭhayogavān\!॥55\!॥\\%linebreak
bindośca stambhane dārḍhyaṃ tadā bhavati yoginaḥ\!।\\%linebreak
tadā ṣoḍaśamabhyāsaṃ nārībhiścinuyādatha\endnote{cinuyād~] M(pc), canuyād M(ac).}\!॥56\!॥\\%linebreak
striyamekaikāṃ dhṛtvā tataścānyāṃ ca\endnote{cānyāṃ ca~] \textit{conj.}, cānyaṃ+ca M.} muhūrte tāṃ tyaktvānyāṃ anyamuhūrte ca, evaṃ ṣoḍaśastrīṇāṃ ratiṃ kṛtvā dine dine, divase kadāpi na kuryāddhīro bindupātaratakrīḍāṃ ca\!। iti mudrāsiddhau\endnote{°siddhau~] \textit{corr.}, °sidhau M.} balamāyuṣpuṣṭivṛddhiśca\endnote{āyuṣpuṣṭi°~] \textit{emend}., āyupuṣṭi° M(pc), āyuṣuṣṭi° M(ac).}\!। prabhavanti yogināṃ vai teja ojo brahmacaryam\!।\\%linebreak
iti koraṇṭakapaddhatau vajrolīmudrā\!।\\%linebreak
\end{sans}
\medskip

\begin{multicols}{2}
\begin{eng}
\begin{enumerate}
\item[48.] Counting [the repetitions] of the practice should be done with the names [of Vāsudeva], starting
from Keśava. It should not be done using Sanskrit or Prakrit numbers, starting from one. Merit
arises from using the names; using numbers is futile.
\end{enumerate}

\heading{[The Pigeon Throat, Frog Belly and Cup Techniques]}

\begin{enumerate}
\item[49.] One should insert the tube into the opening of the penis, blow strongly through the mouth [into
the tube], and hold the breath [within the tube]. This is called “holding the breath” (\tit{vāyor
grahaṇam}).
\item[50.] One should contract the perineal seam below the scrotum, release it, and move [the breath]
forcefully. [If ] one moves [the breath] repeatedly or †performs [contractions] repeatedly, one
acquires the strength of a monitor lizard.†
\item[51.] One should take in air through the tube, form a scissor [shape] with the forefinger and thumb, hold
a gem [with it], and draw [the gem] into the penis with a \tit{huṃ huṃ} sound.
\item[52.] One should expel it with the sound \tit{cārya}, and then with the sound \tit{phūt}. After making the \tit{huṃ
huṃ} sound, one then [makes] the \tit{phūt} sound along with the \tit{cārya} sound.
\item[53.] As is well known, Gahana taught the \tbf{Pigeon Throat Practice}. [Then] the yogi should gradually
perform the \tbf{Frog Belly Practice}.
\item[54.] One should take in air through the tube, fill the bladder using the proper method, and make it like
the belly of a frog. This is called the \tbf{Frog Belly [Practice]}.
\item[55.] After [performing] the \tbf{Frog Belly}, the practice of the \tbf{Cup} is taught. When drawn through the
penis, the air [remains] near the tube. The Haṭhayogī should draw it up through his navel.
\item[56.] And when [the downward flow of ] semen is stopped, the yogi becomes strong. At that time, he
should then acquire [the fruits of ] practising with sixteen women.
\end{enumerate}

\heading{[Afterword]}

The wise [yogi] should embrace one woman and then another, releasing one after a brief period, and
[embracing] the next for another brief period. In this way, he has sex with sixteen women everyday, and
should never engage in amorous play that results in the loss of semen. When the \tit{mudrā} is perfected
thus, strength and the increase of nourishment and longevity occur, and the radiance, vigour, and
celibacy of yogis prevail.

Thus, \tit{vajrolīmudrā} in Koraṇṭaka’s manual.
\end{eng}
\end{multicols}
\bigskip


%~ \begin{multicols}{2}
%~ \begin{kan}
%~ \end{kan}
%~ \end{multicols}
